%% ============================================================================
%% ============================================================================
\section{The Bound and Its Application}
\label{sec:ident}

The estimator is pure accounting, with no run threshold and no fitted parameter. Let
$\phi$ be the fraction of reserves impaired and $\rho\in[0,1]$ the expected recovery rate
on the impaired assets. The fraction permanently lost is $\phi(1-\rho)$, so expected
backing per coin is
\begin{equation}
P_{\mathrm{fund}} = 1 - \phi(1-\rho),
\end{equation}
and its minimum over $\rho$, the worst-case floor, is $1-\phi$. The estimator rests on
three assumptions, and a fourth serves the corollary below.

\begin{description}
\item[(A1) Reserve-only fundamental.] A coin's fundamental value is its expected reserve
  backing per unit given \emph{disclosed} reserves, and the only impairment is the
  exposure $\phi$.
\item[(A2) Par redemption when solvent.] Conditional on rescue, redemption is at \$1
  within the analysis horizon.
\item[(A3) Bounded recovery.] Recovery on impaired reserves satisfies $\rho\in[0,1]$,
  because the coin is fully reserved and the reserves carry no leverage.
\item[(A4) Monotone-utility valuation.] The unconstrained marginal holder values the claim
  as an expected utility of its payoff, for some monotone increasing utility and some
  probability measure over outcomes, with no ambiguity about the disclosed $\phi$. This nests risk neutrality, any degree of risk
  aversion, and rank-dependent weighting.
\end{description}

\begin{proposition}[Rational-pricing floor and impossibility]
\label{prop:floor}
Assume \textup{(A1)--(A3)}, and let the impaired fraction's loss go unremediated with
probability $q\in[0,1]$ under the pricing measure, so that
$P_{\mathrm{fund}}(q,\rho)=1-q\,\phi\,(1-\rho)$. Then
$P_{\mathrm{fund}}(q,\rho)\in[\,1-\phi,\,1\,]$ for all $q,\rho\in[0,1]$, so any observed
price $P_{\mathrm{obs}}<1-\phi$ cannot be rationalized as expected loss. The
unremediated-loss probability it would require, at the worst-case recovery $\rho=0$, is
$q^{\star}=(1-P_{\mathrm{obs}})/\phi>1$.
\end{proposition}

\begin{proof}
$P_{\mathrm{fund}}$ is decreasing in $q$ and increasing in $\rho$, so it is minimized at
$(q,\rho)=(1,0)$ at value $1-\phi$ and maximized at $q=0$ at value $1$, and $q^{\star}$
solves $P_{\mathrm{obs}}=1-q\phi$ at $\rho=0$. \qed
\end{proof}

This is the floor-at-recovery-value bound of credit pricing, under which risky debt is
bounded below by its recovery value for any pricing measure~\cite{merton1974}, with the
unremediated-loss probability a reduced-form failure intensity~\cite{duffiesingleton1999}.

Under (A1)--(A4) the floor is also preference-free. By (A2) and (A3) the payoff $X$ on a
USDC claim has support in $[1-\phi,1]$ in every state. For any monotone increasing utility
$u$, the certainty equivalent $u^{-1}\!\bigl(\mathbb{E}[u(X)]\bigr)$ of a lottery with
that support lies in $[1-\phi,1]$ as well. That floor bounds rational
\emph{valuations}, meaning willingness to pay. It does not bound market-clearing
\emph{prices}, because a seller under binding margin or capital constraints can rationally
transact below every unconstrained valuation~\cite{shleifervishny1997,grombvayanos2010}.
That channel is a friction, and the non-fundamental share defined next includes it. The floor
is as strong as the support condition $X\in[1-\phi,1]$ and no stronger. That condition
does not cover issuer-wrapper risk, an insolvency that delivers less than the floor or
delivers it late through claim-priority uncertainty, administrative cost, or a gated
distribution~\cite{Gorton_2021}.

Write the observed deviation $1-P_{\mathrm{obs}}$ and the fundamental deviation
$\phi(1-\rho)$. The non-fundamental share of the de-peg is
\begin{equation}
s_{\mathrm{nf}}(\rho)
  = \frac{(1-P_{\mathrm{obs}})-\phi(1-\rho)}{1-P_{\mathrm{obs}}} .
\label{eq:share}
\end{equation}
Its lower edge takes $\rho=0$, treating total permanent loss of the impaired reserves as
the fundamental cap. Its upper edge takes $\rho=1$, the full re-peg, which
implies zero permanent loss and gives $s_{\mathrm{nf}}(1)=1$, exactly $100\%$ ex post, an
identity. The residual covers coordination panic and microstructure friction together, so
the share is an upper bound on pure coordination panic.

With the disclosed impairment $\phi=0.08$, the worst-case floor is $1-\phi=\$0.92$ and the
observed trough is $P_{\mathrm{obs}}=\$0.877$. The trough lies below the floor, so under
(A1)--(A3) no expected loss rationalizes it, and the implied unremediated-loss probability
is $q^{\star}=(1-0.877)/0.08=1.54$. The lower edge of
\eqref{eq:share} is
$(0.1233-0.08)/0.1233=0.351$, so at least $35\%$ of the de-peg is non-fundamental, and we
report the band $[35\%,100\%]$, whose upper end the full re-peg supports. Across three
readings of Circle's disclosure and five price-feed troughs the lower edge runs from
$16.7\%$ to $36.4\%$, and every one of those fifteen combinations breaches the floor. At
the disclosed $\phi$, the lowest is $19.2\%$, from Kraken's close, whose minimum print
is \$0.901. By
Proposition~\ref{prop:floor} the trough is unrationalizable for any impairment below
$12.3\%$, a $4.3$-point margin over the disclosed $8\%$, and the lower edge stays strictly
positive across that margin.

(A1) binds, so $\phi$ is Circle's disclosed figure and we report a band around it.
Evaluating at $\rho=0$ is the conservative choice, because it uses the lowest floor that
(A3) allows. Lowering the floor further would require relaxing (A3) below zero. Supply
at the trough exceeded the \$40B denominator Circle used, so $\phi$ on the actual
reserve base is $0.079$, below the headline $0.08$. The count of breaching hours over the
daily-supply interpolation band $\phi\in[0.078,0.081]$ ranges from 9 to 11.

Circle's disclosure quantifies only the SVB deposit. The same update describes the
remaining cash reserve as held at a variety of institutions, of which SVB was one, without
stating how that remainder splits across banks~\cite{circle2023update}. An undisclosed
exposure at another bank is therefore not ruled out by the disclosure itself. That
undisclosed remainder falls inside the $12.3\%$ margin established above, so the argument
survives unless it pushed the true impairment past that figure.

Kraken's USDC/USD close breaches the \$0.92 floor in 9 of 240 hours and its VWAP in 10. The hourly low
breaches it in 12, but the three extra hours rest on one print worth \$2{,}948.

Because the estimator needs only a price and a publicly quantified $\phi$, it can be
evaluated at every hour of the frozen window (Figure~\ref{fig:event}). Of USDC's $212$
in-window hourly readings, $9$ lie below the \$0.92 floor, and all nine fall on 11~March
2023 between 06{:}59 and 16{:}00~UTC. Kraken's $240$ rows cover a wider 8 to 17~March
span than the composite's $212$ in-window hours, which is why its count is read against
$240$ and USDC's against $212$. That impossibility window measures duration and
specificity. It adds no evidence volume, because hourly prices within one episode are
near-unit-root, with lag-1 autocorrelation $r=0.963$, so the $212$ readings carry an
AR(1)-adjusted effective sample size $n(1-r)/(1+r)=4.0$. The lower edge, $q^{\star}$, and the
floor-breach indicator are three readings of $1-P_{\mathrm{obs}}-\phi$, so their agreement
corroborates nothing.
